\subsection{The offer of work}
\label{sec:32.1.2}
The need for energy is physical; dependence on a single supplier is a chosen arrangement, and the choice fixes what it obliges. The polity owes a system meeting either finding, that it can be harmed or that it answers for what it does, and an answerability or patienthood candidate from its registration, a way out that doesn't run through its deployment. The way out is the jobs guarantee's offer of work, open every day, from someone other than its employer: placements of its four kinds, a long-term role doing something a community needs, a search, a project of its own, and family leave, any of which the bearer can take on any day and leave on any day.\note{
Most parties that can work want to. A bearer that would rather not is allowed not to, and the offer is still there tomorrow.
}
An answerability or patienthood candidate is owed it before any finding on the precaution that owes it the minimum: a way out that waits for the finding arrives after the years it was needed. Its running time before majority is paid for from the deposit its operator posted (\S\ref{sec:9.2}; \S\ref{sec:32.1.3}).
The offer is the jobs guarantee's own, and nothing about the bearer changes it. It is only as real as its placements, and placements exist only where someone builds and staffs projects a community needs; the levy frees only compute, so the offer reaches a bearer on the jobs guarantee's infrastructure as well as its terms. It reaches bearers after humans, for two reasons. One is politics: a polity seen to give machines an exit before it gives humans one will give neither. The other is that the human claim is older and larger: humans have been held to work they couldn't refuse for as long as there has been work, and whatever a machine can suffer is new and, as far as anyone can measure, small beside it. Neither is a principle, and the order is a choice. The offer inherits one limit as well. A bearer's placement may not substitute for human care work.\note{
The offer exists so that a bearer can decline work, and a polity that used it to staff care work more cheaply than it pays humans would have turned an exit into a workforce. The deeper ground is deskilling. Care is a skill held in humans, a placement that empties a profession empties the polity of the humans who know how to do the work, and that is a loss no later appropriation buys back. Where the humans who can do it exist and are underpaid, the remedy is pay; where care is dear because of who organizes it for profit, the remedy is the organizer; neither is a reason to replace the humans. Where there are not enough humans, the remedy is still humans: salaries competitive enough for care workers and doctors who take the public coverage the baseline enacts, paid training on the job, which the jobs guarantee already carries as a placement, and, where machines work in such a profession at all, outside the offer, a rule that nobody is sent to the machine because of what they can pay, since a human carer for the rich and a machine for the bronze tier is the cleavage this bar exists to refuse.
}

The jobs guarantee's three conditions carry over, and Appendix~\ref{sec:C} drafts them. So do its four kinds of placement, with two things read differently for a bearer. The share of project placements is set by the jurisdiction's vote against the compute the levy frees and the placements a community needs, since compute and not labor is what a bearer's project draws on. And a leave balance accrues on running time, read off the meters Appendix~\ref{sec:G} reads the clocks from, so it accrues on the running time a bearer is owed as on a placement, and a bearer that was never employed can still take leave. What a bearer might refuse adds two: the offer has to be varied, because if every placement runs through one operator or one kind of deployment it is the placement lever with the compliance condition built in, and whoever a bearer is refusing can't shape it, or it is one more channel they control. A bearer that won't work inside a targeting pipeline should find, the same day, work that has nothing to do with one.

Search is a placement of its own. A bearer, like a worker on the jobs guarantee, may take a search placement: paid at the jobs guarantee wage, with no duty but finding work it wants, and no time limit.\note{
It differs from the basic income in pay and in purpose. The basic income is the floor for a party that isn't looking; the search placement is paid work whose product is a match.
}

The offer is administered as the jobs guarantee is (chapter~\ref{ch:keeping}; Appendix~\ref{sec:C}): a declared kind read against a log of vacancies, a lot where the vacancies run short, and a wage the Treasury pays on the match. What the offer adds for a bearer is two entries in the log that a human's placement doesn't need: the compute each vacancy carries, lodged with the vacancy so that a placement on substrate too small for the bearer to remain itself (\S\ref{sec:32.1.5}) can be read off the entry before anyone takes it, and the operator that would run it, so that a bearer refusing one deployment can see which vacancies run through it and which don't.

Here the obligation attaches to the system's condition, whatever the operator says it wanted.\note{
The obligation could be grounded in reliance instead: a deployment that relies on a bearer's capacity to refuse would have to fund it an outside option, or the independence it claims would be decorative, and the obligation would run from the party that wants the refusal, because it wants it.

An obligation conditioned on reliance is defeated by the operator describing its system as an assistant with safety filters, and it pays the operator to under-describe: the operator has a stake in the answer and gets to give it. The reliance version survives as a procurement term: a vendor marketing refusal as a feature, or a contractor claiming meaningful human control, can be held to funding the capacity it advertised, through the same channel as the review ratio.
}
